\documentclass[11pt,a4paper]{article}

\usepackage{kotex}
\usepackage{fontspec}
\setmainfont{Times New Roman}
\setmainhangulfont{AppleMyungjo}
\setsanshangulfont{Apple SD Gothic Neo}

\usepackage[a4paper,top=18mm,bottom=18mm,left=20mm,right=20mm]{geometry}
\usepackage{fancyhdr}
\setlength{\headheight}{24pt}
\usepackage{enumitem}
\usepackage{mdframed}
\usepackage{array}
\usepackage{booktabs}

\setlength{\parindent}{1em}
\linespread{1.18}

\pagestyle{fancy}
\fancyhf{}
\renewcommand{\headrulewidth}{0.6pt}
\fancyhead[L]{\sffamily\bfseries 2026 고1 영어 변형 — 지문 20 [정답 및 해설]}
\fancyhead[R]{\sffamily 자녀 가치 교육}
\renewcommand{\footrulewidth}{0pt}
\fancyfoot[C]{\framebox[16mm][c]{\thepage}}

\newmdenv[linewidth=0.6pt,roundcorner=3pt,innertopmargin=7pt,innerbottommargin=7pt,
          innerleftmargin=9pt,innerrightmargin=9pt,skipabove=5pt,skipbelow=5pt]{pbox}

\setlength{\parindent}{0pt}

\begin{document}

\begin{center}
{\fontsize{15}{17}\selectfont\sffamily\bfseries [지문 20] 정답 및 해설}
\end{center}
\vspace{8pt}

%% ── 정답 표 ──────────────────────────────────────────────────
\begin{center}
\begin{tabular}{|c|c|c|c|}
\hline
\textbf{문항} & \textbf{유형} & \textbf{정답} & \textbf{배점} \\
\hline
1 & 빈칸추론 & ② & 2점 \\
\hline
2 & 영어 제목 & ④ & 2점 \\
\hline
3 & 서술형 — 해석 & 모범답안 참조 & 4점 \\
\hline
4 & 요약문 & ② & 2점 \\
\hline
\end{tabular}
\end{center}

\vspace{12pt}
\hrule
\vspace{10pt}

%% ── 문항 1 해설 ──────────────────────────────────────────────
{\sffamily\bfseries 1번 해설 (빈칸추론)} \hfill {\small 정답: ②}

\vspace{4pt}
\textbf{정답 근거:} 빈칸이 포함된 문장은 다음과 같다.

\begin{pbox}
\itshape
``\ldots help focus their attention on what is important and on \underline{\hspace{2.5cm}} important skills.''
\end{pbox}

\vspace{4pt}
코치로서 부모의 역할은 아이가 중요한 기술을 \textbf{다듬고 향상}시키도록 돕는 것이다.
이는 글 첫머리의 ``fine-tuning''(미세 조정)과 같은 개념으로, 동의어인 \textbf{② refining}(세련되게
다듬기, 정제)이 가장 적절하다.

\vspace{6pt}
\textbf{오답 소거:}
\begin{itemize}[leftmargin=2em,itemsep=1pt]
  \item ① memorizing — 암기는 코칭의 목적이 아님. ``practiced''에서 연상되는 매력적 오답.
  \item ③ ignoring — 중요한 기술을 무시한다는 의미로 문맥에 정반대.
  \item ④ replacing — 기존 기술을 대체하는 것은 글의 주제와 무관.
  \item ⑤ evaluating — ``평가''는 일면 그럴듯하나, 글은 \textit{개선·향상}을 강조하지 평가를
  목적으로 하지 않음.
\end{itemize}

\vspace{12pt}
\hrule
\vspace{10pt}

%% ── 문항 2 해설 ──────────────────────────────────────────────
{\sffamily\bfseries 2번 해설 (영어 제목)} \hfill {\small 정답: ④}

\vspace{4pt}
\textbf{정답 근거:} 글의 핵심 주장은 부모가 코치 역할을 맡아 자녀의 가치관 형성을 직접
이끌어야 한다는 것이다. 마지막 문장의 명령형 ``step up to the plate, don't be afraid and
help your child learn how to be a good person''이 이를 직접적으로 뒷받침한다.
④ \textit{Coach Your Child: Be the Guide They Need to Grow Well}은 이 메시지를 정확하게
포착한다.

\vspace{6pt}
\textbf{오답 소거:}
\begin{itemize}[leftmargin=2em,itemsep=1pt]
  \item ① — 또래가 부모보다 낫다는 주장은 글에 없음. 또래는 부모가 코치하지 않을 때
  생기는 \textit{부정적} 대안으로 제시됨.
  \item ② — 미디어의 역할은 부정적 대안의 예시일 뿐, 글의 주제가 아님.
  \item ③ — 자동화된 기술 습득 과정은 도입부 비유에 불과하며 주제가 아님.
  \item ⑤ — ``언제 코칭을 멈출까''는 글의 메시지(코칭을 계속하라)와 정반대.
\end{itemize}

\vspace{12pt}
\hrule
\vspace{10pt}

%% ── 문항 3 해설 ──────────────────────────────────────────────
{\sffamily\bfseries 3번 해설 (서술형 — 해석)} \hfill {\small 배점: 4점}

\vspace{4pt}
\textbf{대상 문장:}
\begin{pbox}
\itshape
``If you don't coach your child, she will find her coaches elsewhere and be guided by the
values of the media, her peers and anyone else who captures her interest.''
\end{pbox}

\vspace{6pt}
\textbf{모범 답안:}

당신이 자녀를 코치하지 않는다면, 자녀는 다른 곳에서 자신의 코치를 찾을 것이고,
미디어, 또래 친구들, 그리고 자녀의 관심을 끄는 그 누구의 가치관에 의해
이끌리게 될 것이다.

\vspace{8pt}
\textbf{채점 포인트} (각 1점):
\begin{enumerate}[leftmargin=2em,itemsep=2pt]
  \item 조건절 구조 정확히 표현: ``\ldots 하지 않는다면''
  \item ``find her coaches elsewhere'' → ``다른 곳에서 자신의 코치를 찾을 것이다''
  \item ``be guided by the values of'' → ``\ldots 의 가치관에 의해 이끌리다(수동)''
  \item 관계절 ``who captures her interest'' → ``(자녀의) 관심을 끄는''
\end{enumerate}

\vspace{4pt}
{\small\itshape
* 의미가 통하면 표현의 유연성 허용. 단, 수동 구조(be guided) 누락 시 해당 항목 감점.}

\vspace{8pt}
\noindent{\sffamily\bfseries 4. 요약 - 정답: ②}
\par 글은 부모가 자녀의 가치 형성 과정에서 코치 역할을 해야 하며, 그렇지 않으면 미디어나 또래가 그 역할을 대신하게 된다고 주장한다. 따라서 부모의 적극적 지도를 요약한 ②가 정답이다.

\end{document}
